\begin{theorem}[Cauchy–Schwarz]
For real numbers \(a_{1},\dots ,a_{n}\) and \(b_{1},\dots ,b_{n}\),
\[
\left(\sum_{i=1}^{n}a_{i}b_{i}\right)^{2}
\le
\left(\sum_{i=1}^{n}a_{i}^{2}\right)
\left(\sum_{i=1}^{n}b_{i}^{2}\right).
\]
Equality holds iff there exists \(\lambda\in\mathbb{R}\) such that
\(a_{i}=\lambda b_{i}\) for every \(i\) (including the trivial case
\(b_{1}=\dots =b_{n}=0\) or \(a_{1}=\dots =a_{n}=0\)).
\end{theorem}

\begin{proof}
\textbf{Degenerate case \(\displaystyle\sum_{i=1}^{n}b_{i}^{2}=0\).}
If \(\sum_{i=1}^{n}b_{i}^{2}=0\) then each \(b_{i}=0\). Consequently
\[
\sum_{i=1}^{n}a_{i}b_{i}=0,\qquad
\bigl(\sum_{i=1}^{n}a_{i}b_{i}\bigr)^{2}=0,\qquad
\bigl(\sum_{i=1}^{n}a_{i}^{2}\bigr)
\bigl(\sum_{i=1}^{n}b_{i}^{2}\bigr)=0,
\]
so the inequality holds with equality.  
From now on assume \(\displaystyle\sum_{i=1}^{n}b_{i}^{2}>0\).

\medskip
\textbf{Construction of a non‑negative quadratic.}
For an arbitrary real parameter \(t\) define
\[
Q(t):=\sum_{i=1}^{n}\bigl(a_{i}-t\,b_{i}\bigr)^{2}.
\]
Each summand is a square, hence \(Q(t)\ge0\) for all \(t\in\mathbb{R}\).

Expanding and collecting powers of \(t\) gives
\begin{align*}
Q(t)&=\sum_{i=1}^{n}\bigl(a_{i}^{2}-2t\,a_{i}b_{i}+t^{2}b_{i}^{2}\bigr)\\
    &=\underbrace{\sum_{i=1}^{n}b_{i}^{2}}_{A}\,t^{2}
      +\underbrace{\bigl(-2\sum_{i=1}^{n}a_{i}b_{i}\bigr)}_{B}\,t
      +\underbrace{\sum_{i=1}^{n}a_{i}^{2}}_{C}.
\end{align*}
Thus \(Q(t)=At^{2}+Bt+C\) with
\[
A=\sum_{i=1}^{n}b_{i}^{2}>0,\qquad
B=-2\sum_{i=1}^{n}a_{i}b_{i},\qquad
C=\sum_{i=1}^{n}a_{i}^{2}.
\]

\medskip
\textbf{Discriminant condition for a non‑negative quadratic.}
\emph{Lemma.} If a quadratic \(f(t)=At^{2}+Bt+C\) satisfies \(f(t)\ge0\) for
all \(t\in\mathbb{R}\) and \(A>0\), then its discriminant is non‑positive:
\(B^{2}-4AC\le0\).

\emph{Proof of Lemma.} If \(B^{2}-4AC>0\) the polynomial would have two
distinct real roots and would be negative between them, contradicting
\(f(t)\ge0\).  If \(B^{2}-4AC=0\) the polynomial has a double root and is
non‑negative; if \(B^{2}-4AC<0\) it has no real root and is strictly
positive.  ∎

Applying the lemma to \(Q(t)\) yields
\[
B^{2}-4AC\le0.
\]
Substituting the expressions for \(A,B,C\) we obtain
\[
\bigl(-2\sum_{i=1}^{n}a_{i}b_{i}\bigr)^{2}
-4\Bigl(\sum_{i=1}^{n}b_{i}^{2}\Bigr)
   \Bigl(\sum_{i=1}^{n}a_{i}^{2}\Bigr)\le0.
\]

\medskip
\textbf{Algebraic simplification.}
Using the elementary identity
\[
\bigl(-2\sum a_{i}b_{i}\bigr)^{2}
-4\Bigl(\sum b_{i}^{2}\Bigr)\Bigl(\sum a_{i}^{2}\Bigr)
=4\Bigl[\bigl(\sum a_{i}b_{i}\bigr)^{2}
      -\bigl(\sum a_{i}^{2}\bigr)\bigl(\sum b_{i}^{2}\bigr)\Bigr],
\]
the previous inequality becomes
\[
4\Bigl[\bigl(\sum a_{i}b_{i}\bigr)^{2}
      -\bigl(\sum a_{i}^{2}\bigr)\bigl(\sum b_{i}^{2}\bigr)\Bigr]\le0.
\]
Dividing by the positive constant \(4\) gives the Cauchy–Schwarz inequality
\[
\boxed{\displaystyle
\bigl(\sum_{i=1}^{n}a_{i}b_{i}\bigr)^{2}
\le
\bigl(\sum_{i=1}^{n}a_{i}^{2}\bigr)
\bigl(\sum_{i=1}^{n}b_{i}^{2}\bigr)}.
\]

\medskip
\textbf{Equality condition.}
Equality in the discriminant inequality occurs exactly when
\(B^{2}-4AC=0\), i.e. when the quadratic \(Q(t)\) has a double root.
Hence there exists \(t_{0}\in\mathbb{R}\) such that \(Q(t_{0})=0\).
Since
\[
Q(t_{0})=\sum_{i=1}^{n}\bigl(a_{i}-t_{0}b_{i}\bigr)^{2}=0,
\]
each summand must be zero, giving
\[
a_{i}=t_{0}\,b_{i}\qquad\text{for all }i=1,\dots ,n.
\]
Conversely, if a constant \(\lambda\) satisfies \(a_{i}=\lambda b_{i}\) for
all \(i\), then
\[
\bigl(\sum a_{i}b_{i}\bigr)^{2}
=\bigl(\lambda\sum b_{i}^{2}\bigr)^{2}
=\lambda^{2}\bigl(\sum b_{i}^{2}\bigr)^{2}
=\bigl(\sum a_{i}^{2}\bigr)\bigl(\sum b_{i}^{2}\bigr),
\]
so equality holds.  The case \(\lambda=0\) corresponds to the
degenerate situation where one of the vectors is the zero vector.

Thus equality holds iff there exists \(\lambda\in\mathbb{R}\) such that
\(a_{i}=\lambda b_{i}\) for every index \(i\) (including the trivial
zero‑vector case). \(\square\)
\end{proof}